% Generated draft: requires mathematical and source review.
\documentclass[11pt]{article}
\usepackage[margin=25mm]{geometry}
\usepackage{fontspec}
\setmainfont{Latin Modern Roman}
\usepackage{amsmath,amssymb,mathtools,unicode-math}
\setmathfont{Latin Modern Math}
\usepackage{enumitem,needspace,xurl,hyperref}
\setlength{\parindent}{0pt}
\setlength{\parskip}{5pt}
\setlength{\emergencystretch}{3em}
\newcommand{\N}{\mathbb{N}}
\newcommand{\Z}{\mathbb{Z}}
\newcommand{\Q}{\mathbb{Q}}
\newcommand{\R}{\mathbb{R}}
\newcommand{\C}{\mathbb{C}}
\newcommand{\F}{\mathbb{F}}
\newcommand{\A}{\mathbb{A}}
\newcommand{\PP}{\mathbb{P}}
\newcommand{\E}{\mathbb{E}}
\newcommand{\eps}{\varepsilon}
\newcommand{\dd}{\,\mathrm{d}}
\begin{document}
\section*{UnsolvedMath ID: 30005189}
\textbf{Source problem code:} catalogue entry
\subsection{Definitions and mathematical statement}
The supplied record does not specify the domain or hypotheses for \(k\), the meaning of \(\operatorname{Var}_k\), the construction of its groups \(K_0\) and \(K_1\), the meaning of \(S\) and \(\pi_1S\), the tensor-product convention, or the construction of the displayed map. Retaining the supplied notation, the question is whether the intended map
\[
\varphi\colon K_0(\operatorname{Var}_k)\otimes\pi_1S
\longrightarrow K_1(\operatorname{Var}_k)
\]
is surjective. Here surjectivity means
\[
\forall y\in K_1(\operatorname{Var}_k)\quad
\exists x\in K_0(\operatorname{Var}_k)\otimes\pi_1S
\quad \varphi(x)=y.
\]
This quantifier formulation is conditional on specifying the intended objects and map; it does not resolve their identities or supply missing hypotheses on \(k\).

The accompanying passage states that an affirmative answer to a second question would imply an affirmative answer to a first question, and therefore focuses on the second question. The first question and sufficient surrounding context are absent, so this implication cannot be formulated precisely from the supplied record.
\subsection{Short English statement}
Does every element of \(K_1(\operatorname{Var}_k)\) lie in the image of the displayed map? The intended map and several underlying conventions remain unresolved in the supplied record, so this is a conditional transcription requiring mathematical review.
\subsection{Sources}
\begin{itemize}
\item Supplied local problem registry: permanent UnsolvedMath ID \(30005189\), source code \texttt{catalogue entry}, title ``Surjectivity in Algebraic K-Theory of Varieties.'' No individual proposer is supplied; the recorded proposed year is 2022.
\item Supplied report citation: \emph{Non-Commutative Geometry and Cyclic Homology} (2023), Oberwolfach Reports, pp.~1987--2058, DOI \texttt{10.4171/owr/2022/35}. Source URL and supplied literature source labelled ``original Oberwolfach report'': \url{https://doi.org/10.4171/owr/2022/35}. The record attributes the problem to workshop open-problem discussions documented in this report.
\item Literature assessment dated exactly \texttt{2026-08-22}: the supplied curation labels the problem open, but states that the coefficient notation and preceding question are corrupted, preventing reliable identification of the precise \(K\)-theory map and assignment of a definitive surjectivity result to this exact statement. This is supplied dated curation, not an independently verified literature assessment.
\item The registry also labels the statement \texttt{corrected\_verified} and reports that it was checked against the official report and preserves its hypotheses, notation, and target. That supplied verification claim conflicts with the unresolved formulation issues recorded in its literature assessment; no source visit or independent verification is asserted here.
\end{itemize}
\end{document}
